% Options for packages loaded elsewhere
\PassOptionsToPackage{unicode}{hyperref}
\PassOptionsToPackage{hyphens}{url}
%
\documentclass[
  11pt,
]{article}
\usepackage{amsmath,amssymb}
\usepackage{iftex}
\ifPDFTeX
  \usepackage[T1]{fontenc}
  \usepackage[utf8]{inputenc}
  \usepackage{textcomp} % provide euro and other symbols
\else % if luatex or xetex
  \usepackage{unicode-math} % this also loads fontspec
  \defaultfontfeatures{Scale=MatchLowercase}
  \defaultfontfeatures[\rmfamily]{Ligatures=TeX,Scale=1}
\fi
\usepackage{lmodern}
\ifPDFTeX\else
  % xetex/luatex font selection
\fi
% Use upquote if available, for straight quotes in verbatim environments
\IfFileExists{upquote.sty}{\usepackage{upquote}}{}
\IfFileExists{microtype.sty}{% use microtype if available
  \usepackage[]{microtype}
  \UseMicrotypeSet[protrusion]{basicmath} % disable protrusion for tt fonts
}{}
\makeatletter
\@ifundefined{KOMAClassName}{% if non-KOMA class
  \IfFileExists{parskip.sty}{%
    \usepackage{parskip}
  }{% else
    \setlength{\parindent}{0pt}
    \setlength{\parskip}{6pt plus 2pt minus 1pt}}
}{% if KOMA class
  \KOMAoptions{parskip=half}}
\makeatother
\usepackage{xcolor}
\usepackage[margin=25mm]{geometry}
\setlength{\emergencystretch}{3em} % prevent overfull lines
\providecommand{\tightlist}{%
  \setlength{\itemsep}{0pt}\setlength{\parskip}{0pt}}
\setcounter{secnumdepth}{-\maxdimen} % remove section numbering
\usepackage{amssymb,amsmath}
\ifLuaTeX
  \usepackage{selnolig}  % disable illegal ligatures
\fi
\usepackage{bookmark}
\IfFileExists{xurl.sty}{\usepackage{xurl}}{} % add URL line breaks if available
\urlstyle{same}
\hypersetup{
  pdftitle={Equal Gaussian cells: the sharp simplex first-moment theorem},
  pdfauthor={mxym/math research project},
  hidelinks,
  pdfcreator={LaTeX via pandoc}}

\title{Equal Gaussian cells: the sharp simplex first-moment theorem}
\author{mxym/math research project}
\date{8 October 2026}

\begin{document}
\maketitle

\section{Abstract}\label{abstract}

For every number of cells \(k\ge2\), we determine the sharp squared
first-moment sum for equal-probability Gaussian partitions when the
ambient dimension is at least \(k-1\). The unique maximizers are central
regular simplex cones, and the value is
\((\mathbb E\max_{i\le k}Z_i)^2/(k-1)\) for independent standard
normals. The bound also holds in lower dimensions and for fractional
partitions. We prove a covariance comparison by combining Gaussian flux
with the established Gaussian multi-bubble perimeter theorem. A weighted
Cauchy inequality yields a first-order differential inequality along
each path from the regular covariance; its initial derivative vanishes
by symmetry. This gives the global bound and equality classification,
with an exact nonnegative integral expression for the deficit. No
numerical search, boundary variation, local Hessian or mountain-pass
argument is needed. We distinguish the imported perimeter theorem and
partial formal checks from the new analytic comparison.

\section{1. Statement and prior
conjecture}\label{statement-and-prior-conjecture}

Let \(\gamma_d\) be standard Gaussian probability on \(\mathbb R^d\).
For a measurable partition \(C_1,\ldots,C_k\) define
\(m_i=\int_{C_i}x\,d\gamma_d\) and \(F=\sum_i|m_i|^2\). Write
\(Z_1,\ldots,Z_k\) for independent standard normals and

\[
 a_k=\mathbb E\max_{i\le k}Z_i,
 \qquad c_k=\frac{a_k}{\sqrt{k-1}},\qquad F_k=c_k^2.
\]

These constants are exact; for example integration by parts gives

\[
 a_k=k(k-1)\int_{\mathbb R}\phi(x)^2\Phi(x)^{k-2}\,dx>0,       \tag{1}
\]

where \(\phi,\Phi\) are the standard normal density and distribution.

\textbf{Theorem 1.} For every integer \(k\ge2\), every integer
\(d\ge1\), and every Gaussian partition with \(\gamma_d(C_i)=1/k\),

\[
 \sum_{i=1}^k\left|\int_{C_i}x\,d\gamma_d(x)\right|^2
             \le\frac{a_k^2}{k-1}.                         \tag{2}
\]

If \(d\ge k-1\), equality holds precisely for the central regular
simplex winning partition in a \((k-1)\)-dimensional subspace, with the
orthogonal complement unrestricted, up to null sets, orthogonal
transformations and relabeling. If \(d<k-1\), the inequality is strict
for every partition; we do not claim this is the sharp bound at a fixed
such lower dimension.

The identical statement holds for fractional partitions:
\(0\le f_i\le1\), \(\sum_i f_i=1\) almost everywhere, and
\(\int f_i\,d\gamma_d=1/k\), using \(m_i=\int xf_i\,d\gamma_d\).
Equality forces the same ordinary winning indicators.

This settles the equal-mass subcase for every number of cells of the
first-moment conjecture in {[}H19, Problem 1.15 and Conjecture 1.16{]},
and in particular the four-cell dimension-three question {[}H14,
Conjecture 3{]}. It is not the positive-correlation noise-stability
conjecture or the arbitrary prescribed-mass statement. The latter is
false as written {[}EX{]}. Our \href{LITERATURE_STATUS.md}{literature
status} reports the precise sources checked and makes no worldwide
first-proof claim.

\section{2. The balanced covariance value
function}\label{the-balanced-covariance-value-function}

Put \(\mathbf1=(1,\ldots,1)^T\in\mathbb R^k\),
\(P=I-\mathbf1\mathbf1^T/k\), \(n=k-1\),

\[
 \mathcal K=\{Q=Q^T\succeq0:Q\mathbf1=0,\ \operatorname{tr}Q=1\},
 \qquad Q_*=P/n.
\]

For every positive semidefinite \(Q\) annihilating \(\mathbf1\), define

\[
 \mathcal C(Q)=\inf_{\lambda\in\mathbb R^k}
 \left\{\frac1k\sum_i\lambda_i+
               \mathbb E\max_i(Y_i-\lambda_i)\right\},
 \qquad Y\sim N(0,Q).                                      \tag{3}
\]

For a centered \(k\times n\) score matrix \(M\) with \(MM^T=Q\), this is
called the balanced score value. For distinct rows it equals the optimal
assignment value, as proved below. Prices attain their infimum: in the
gauge \(\min\lambda_i=0\), the expression is at least
\(\max\lambda_i/k\) by using a zero-priced mean-zero score. For distinct
scores \(v_i\), differentiating in prices gives \(1/k-\gamma_n(C_i)\),
where

\[
 C_i=\{x:v_i\cdot x-\lambda_i\ge v_j\cdot x-\lambda_j\ \forall j\}.
                                                               \tag{4}
\]

Ties are Gaussian-null. Minimizing prices therefore give all masses
\(1/k\); the winning labels attain the linear assignment pointwise.

Two elementary facts include the singular covariance boundary:

\[
 \mathcal C(sQ)=\sqrt{s}\,\mathcal C(Q)\ (s\ge0),
 \qquad \mathcal C\text{ is continuous on the closed cone}.       \tag{5}
\]

For homogeneity scale scores and prices, and note \(\mathcal C(0)=0\).
For continuity couple \(Q^{1/2}G\) and \(\widetilde Q^{1/2}G\) on the
same standard Gaussian space. The difference of the maxima, for any
price, is at most \(|(Q^{1/2}-\widetilde Q^{1/2})G|\). Its expectation
bounds the difference of the two infima, uniformly in prices; the matrix
square root is continuous.

We use differentiability only where \(Q\) is positive definite on
\(\mathbf1^\perp\). There the \(k\) scores are affinely independent in
\(\mathbb R^n\). Every pair of winning cells has a positive Gaussian
facet area: the nonsingular score-difference coordinates allow any
chosen pair to tie at the top while all other scores are strictly lower.
For the ordinary \(ij\) facet put

\[
 \ell_{ij}=|v_i-v_j|,
 \quad A_{ij}=\int_{\Sigma_{ij}}\varphi_n\,d\mathcal H^{n-1},
 \quad w_{ij}=A_{ij}/\ell_{ij},
 \quad (Lu)_i=\sum_{j\ne i}w_{ij}(u_i-u_j).               \tag{6}
\]

Let \(B\) have rows \(b_i=\int_{C_i}x\,d\gamma_n\). Gaussian integration
by parts gives

\[
 B=LM,\quad \mathcal C(Q)=\langle M,B\rangle
      =\operatorname{tr}(LQ)=\sum_{i<j}w_{ij}\ell_{ij}^2,
 \quad \operatorname{tr}L=2\sum_{i<j}w_{ij}.             \tag{7}
\]

The inward \(ij\) normal of cell \(i\) is \((v_i-v_j)/\ell_{ij}\). To
justify flux on unbounded cells truncate with large balls; their
additional Gaussian boundary term tends to zero exponentially. The price
term in (3) cancels because all winning masses are \(1/k\).

\textbf{Lemma 2 (smooth derivative).} On this positive cone
\(\mathcal C\) is smooth, and for every symmetric \(D\) annihilating
\(\mathbf1\),

\[
 d\mathcal C(Q)[D]=\tfrac12\operatorname{tr}(LD).         \tag{8}
\]

\emph{Proof.} The \(n\) score differences against label \(k\) give a
nonsingular affine change of coordinates. Its winning regions are fixed
polyhedral cones. Gaussian density derivatives in these coordinates are
locally uniformly dominated by a polynomial times a Gaussian, including
for the linear integrand defining the expected maximum. Thus the score
and price integrals are smooth. The price Hessian is \(L\): moving a
price moves its pairwise facet by that price gap divided by
\(\ell_{ij}\). Since every weight is positive, \(L\) is positive on the
gauge \(\sum\lambda_i=0\). It is positive along any price segment,
giving unique prices in this gauge. The implicit function theorem makes
the balancing prices smooth.

Write \(D=MAM^T\) with \(A=A^T\) and realize the covariance path using
\(M(t)=M(I+tA)^{1/2}\). Every such \(D\) has this representation since
\(M\) maps onto \(\mathbf1^\perp\). The envelope score derivative is
\(B\). At zero, \(M'=MA/2\), so (7) gives
\(d\mathcal C(Q)[D]=\operatorname{tr}(B^TMA)/2
=\operatorname{tr}(LD)/2\). \(\square\)

\section{3. The imported perimeter theorem and the regular
constant}\label{the-imported-perimeter-theorem-and-the-regular-constant}

We use the established Gaussian multi-bubble theorem {[}MN22, Theorem
1.1{]}: for \(2\le k\le n+1\) and fixed positive masses in
\(\mathbb R^n\), a simplicial Voronoi cluster of \(k\) equidistant sites
minimizes the total Gaussian perimeter
\(S=\frac12\sum_i P_\gamma(C_i)=\sum_{i<j}A_{ij}\). We always apply it
in dimension \(n=k-1\). At equal masses this minimizer is the central
regular simplex fan: balancing prices of regular scores are unique
modulo constants by Lemma 2, and permutation symmetry makes all of them
equal. Only the perimeter minimum, not its uniqueness part, is needed
below.

For the trace-one regular scores with Gram matrix \(Q_*=P/n\), we can
write \(Y_i=(Z_i-\overline Z)/\sqrt n\). Their winning labels are
uniform, so

\[
 \mathcal C(Q_*)=c_k.
\]

All facet weights are equal by symmetry; (7) gives \(L_*=c_kP\) and
\(w_{ij}^*=c_k/k\). All pair distances are \(\ell_* =\sqrt{2/n}\).
Consequently the minimum perimeter is

\[
 S_* =\frac{k(k-1)}2\frac{c_k}k\sqrt{2/n}
              =c_k\sqrt{n/2}.                            \tag{9}
\]

For any full-rank \(Q\in\mathcal K\), weighted Cauchy--Schwarz and (7)
give

\[
 S^2=\left(\sum_{i<j}w_{ij}\ell_{ij}\right)^2
 \le\left(\sum_{i<j}w_{ij}\right)
        \left(\sum_{i<j}w_{ij}\ell_{ij}^2\right)
 =\tfrac12\operatorname{tr}L\,\mathcal C(Q).
\]

Combining this with \(S\ge S_*\) proves the fundamental inequality

\[
 \mathcal C(Q)\frac{\operatorname{tr}L}{n}\ge c_k^2.       \tag{10}
\]

All weights here are positive, so equality in weighted Cauchy is
possible only when all pair distances are equal. This observation will
supply the equality classification.

\section{4. The radial differential comparison and its
deficit}\label{the-radial-differential-comparison-and-its-deficit}

\textbf{Theorem 3 (covariance form).} For every \(Q\in\mathcal K\),
\(\mathcal C(Q)\le c_k\), with equality only at \(Q_*\).

\emph{Proof.} Put \(Q_t=Q_*+t(Q-Q_*)\). It is positive on
\(\mathbf1^\perp\) for \(0\le t<1\), even if the endpoint \(Q\) is
singular. Let \(C(t)=\mathcal C(Q_t)\) and \(L_t\) be its facet matrix.
For \(0<t<1\), (7)--(8) give

\[
 tC'(t)=\tfrac12\operatorname{tr}(L_t(Q_t-Q_*))
       =\tfrac12\left(C(t)-\frac{\operatorname{tr}L_t}{n}\right).
                                                               \tag{11}
\]

Let \(h(t)=C(t)^2-c_k^2\). Inequality (10) implies

\[
 th'(t)=C(t)^2-C(t)\frac{\operatorname{tr}L_t}{n}
       \le h(t),\qquad
 \left(\frac{h(t)}t\right)'\le0.                         \tag{12}
\]

At zero \(L_*=c_kP\), and \(\operatorname{tr}(Q-Q_*)=0\). Thus (8) gives
\(C'(0)=0\), while \(C(0)=c_k\). It follows that
\(\lim_{t\downarrow0}h(t)/t=0\). The quotient is nonincreasing by (12),
hence \(h(t)\le0\) on \((0,1)\). Continuity (5) gives the claimed bound
at \(t=1\).

If \(\mathcal C(Q)=c_k\), the quotient tends to zero also at one. A
nonincreasing function with both endpoint limits zero is identically
zero. Therefore \(C(t)=c_k\) for every \(0<t<1\), and (11) gives
\(\operatorname{tr}L_t=nc_k\). The perimeter/Cauchy chain in Section 3
then consists entirely of equalities. All \(\ell_{ij}(t)\) are equal,
since all weights are positive. Centered scores of squared norm sum one
with equal pair distances have Gram matrix \(P/n\): summing their
squared pair distances over \(j\) gives \(k|v_i|^2+1\), independent of
\(i\); thus \(|v_i|^2=1/k\) and \(\ell_{ij}^2=2/n\). Hence \(Q_t=Q_*\).
For any fixed \(0<t<1\), this forces \(Q=Q_*\). \(\square\)

There is also an exact deficit identity. For \(Q_t\) as above put

\[
 D(t)=C(t)\frac{\operatorname{tr}L_t}{n}-c_k^2\ge0.
\]

The equality version of (12) reads \((h/t)'=-D(t)/t^2\), giving

\[
 c_k^2-\mathcal C(Q)^2=\int_0^1\frac{D(t)}{t^2}\,dt.     \tag{13}
\]

This is a nonnegative improper integral. At zero smoothness gives
\(h=O(t^2)\); integrating on compact subintervals and using the endpoint
limits proves convergence and (13), even for a singular endpoint. In
particular, writing \(W(t)=\sum_{i<j}w_{ij}(t)\),

\[
 D(t)=\frac2n\left[
   W(t)\sum_{i<j}w_{ij}(t)\ell_{ij}(t)^2-S(t)^2
                         +S(t)^2-S_*^2\right].           \tag{14}
\]

Both terms are nonnegative: the first is the weighted edge-length
variance, and the second is the established perimeter deficit. This is
an exact identity, not a claim of an explicit metric stability constant
or of full covariance concavity.

\section{5. Return to measurable and fractional
partitions}\label{return-to-measurable-and-fractional-partitions}

For any partition in Theorem 1, let \(M\) be its \(k\times d\) moment
matrix. Its rows sum to zero, so \(MM^T\) has rank at most \(k-1\) and
trace \(F\). For every price, its labels give a feasible assignment in
(3):

\[
 F=\mathbb E[m_I\cdot X]\le\mathcal C(MM^T).
\]

This remains valid for fractional labels by replacing the selected score
with their convex combination. The Gaussian score distribution has the
same covariance even if the original dimension differs from \(k-1\).
When \(F>0\), (5) and Theorem 3 give \(F\le c_k\sqrt F\), proving (2).
For \(F=0\) the positive bound is strict.

Equality forces \(MM^T=F_kP/n\) and equality in the feasible assignment.
The regular centered scores have null ties and equal prices by symmetry
and uniqueness, so the original labels must be their winning labels
almost everywhere. Fractional labels likewise must be the unique winning
indicators. Their span has dimension \(k-1\), which excludes equality if
\(d<k-1\).

Conversely, take the regular scores with Gram matrix \(P/n\) and their
central winning cells. By (7) and \(L_*=c_kP\), their actual moments are
\(c_kv_i\), of squared sum \(c_k^2\). Cylindrical extension preserves
both masses and moments. This proves attainment and all equality
assertions. For \(k=2\) the trace-one covariance domain is a singleton;
the proof includes this case. The excluded \(k=1\) case is separately
trivial, with one cell and zero first moment.

\section{6. Mathematical and verification
dependencies}\label{mathematical-and-verification-dependencies}

The only imported non-elementary geometric result is the Gaussian
multi-bubble perimeter theorem {[}MN22{]}. Its hypotheses, ambient
dimension and perimeter normalization are stated in Section 3. The proof
uses Gaussian integration by parts, dominated convergence, smooth matrix
square roots in the positive cone, the implicit function theorem, and
one-variable differential comparison. No Gaussian scalar isoperimetry,
local simplex Hessian, covariance concavity, boundary differentiation,
rank classification or mountain-pass theorem is assumed.

No finite search, solver, floating integral or numerical eigenvalue is a
premise of this result. Partial Lean proofs and exact arithmetic
controls in the package have their own explicit scope; they do not
formalize Gaussian measure, the imported multi-bubble theorem or the
complete endpoint. Internal model reviews are not external human peer
review. The complete earlier four-cell proof {[}FOUR{]} is preserved as
a historical, longer argument; the present proof is standalone and
covers all \(k\).

\section{References}\label{references}

\begin{itemize}
\tightlist
\item
  {[}H14{]} S. Heilman, \emph{Euclidean Partitions Optimizing Noise
  Stability}, Electronic Journal of Probability 19 (2014), no. 71,
  1--37; arXiv:1211.7138v2, Conjecture 3 and Definition 1.7.
\item
  {[}H19{]} S. Heilman, \emph{Stable Gaussian Minimal Bubbles},
  arXiv:1901.03934v1, Problem 1.15 and Conjecture 1.16.
\item
  {[}MN22{]} E. Milman and J. Neeman, \emph{The Gaussian Double-Bubble
  and Multi-Bubble Conjectures}, Annals of Mathematics 195 (2022), no.
  1, 89--206; DOI 10.4007/annals.2022.195.1.2; arXiv:1805.10961v3,
  Theorem 1.1.
\item
  {[}EX{]} mxym/math research project, \emph{A counterexample to the
  prescribed-mass Gaussian regular-simplex conjecture}, repository
  companion research/gaussian-fixed-mass-propeller-counterexample, 2026.
\item
  {[}FOUR{]} mxym/math research project, \emph{Four equal Gaussian
  cells: covariance deformation and tetrahedral rigidity}, complete
  eleven-page repository companion
  research/gaussian-balanced-four-global, commit 0d27bab, 2026.
\end{itemize}

\end{document}
